% ============================================================
% CAPITOLO 42 - RIEPILOGO DELLA PARTE VIII E DELLE PROGRESSIONI
% ============================================================
\chapter{Riepilogo della Parte VIII e delle progressioni}
\label{chap:riepilogo_viii}

\section{Napoletana e seste eccedenti a confronto}
\label{sec:quadro_napoletana}

\begin{center}
\small
\begin{tabularx}{\linewidth}{lccL}
\toprule
\textbf{Voce / nota (sesta napoletana)} & \textbf{La minore} & \textbf{Do maggiore} & \textbf{Movimento} \\
\midrule
Basso (terza) & Re & Fa & sale di grado alla fondamentale della dominante (Mi / Sol) \\
Nota napoletana (II$\flat$) & Si$\flat$ & Re$\flat$ & scende alla sensibile (Sol$\sharp$ / Si); mai raddoppiata \\
Quinta (VI grado) & Fa & La$\flat$ & scende di grado alla fondamentale della dominante \\
Terza raddoppiata & Re & Fa & resta ferma (settima della \sigla{V}{7}{}) o scende a Si / Re \\
\bottomrule
\end{tabularx}
\end{center}

\begin{center}
\small
\begin{tabular}{lllp{5cm}}
\toprule
\textbf{Sesta eccedente} & \textbf{Note (La minore)} & \textbf{Cifratura} & \textbf{Risoluzione} \\
\midrule
Italiana & Fa -- La -- Re$\sharp$ (La raddoppiato) & $\sharp$6 & Fa $\searrow$ Mi, Re$\sharp$ $\nearrow$ Mi, La $\searrow$ Sol$\sharp$, La $\nearrow$ Si: V completa \\
Francese & Fa -- La -- Si -- Re$\sharp$ & $\sharp$6 4 3 & Fa $\searrow$ Mi, Re$\sharp$ $\nearrow$ Mi, La $\searrow$ Sol$\sharp$, Si ferma: V completa \\
Tedesca & Fa -- La -- Do -- Re$\sharp$ & $\sharp$6 5 3 & Fa $\searrow$ Mi, Re$\sharp$ $\nearrow$ Mi, La e Do fermi: \sigla{i}{6}{4}; poi La $\searrow$ Sol$\sharp$, Do $\searrow$ Si: V \\
\bottomrule
\end{tabular}
\end{center}

La sesta napoletana è trattata nel Capitolo \ref{chap:sesta_napoletana}, le seste eccedenti nel Capitolo \ref{chap:seste_eccedenti}. Entrambe preparano la dominante: nella napoletana il basso sale di grado (dal IV al V), nelle seste eccedenti scende di semitono (dal VI al V) mentre la sesta si espande verso l'ottava sulla dominante.

\needspace{0.6\textheight}
\section{Quando usare quale progressione}
\label{sec:sintesi_progressioni}

\begin{regola}[Guida alla scelta delle progressioni]
	\label{reg:scelta_progressioni}

	\begin{tabular}{p{3.6cm}p{4.2cm}p{2.8cm}p{1.6cm}}
		\toprule
		\textbf{Tipo di basso} & \textbf{Progressione consigliata} & \textbf{Effetto} & \textbf{Cap.} \\
		\midrule
		Salita scalare lunga (5+ note) & Regola dell'Ottava ascendente (alternanza 5, 6/5) & Tensione progressiva verso la cadenza & \ref{chap:ottava_preliminare}, \ref{chap:regola_ottava_completa} \\
		\midrule
		Discesa scalare & Progressione per terze o Regola dell'Ottava & Dolcezza, malinconia & \ref{chap:progressioni_armoniche} \\
		\midrule
		Movimento circolare & Quinte discendenti & Forza, direzione, conclusività & \ref{chap:progressioni_armoniche} \\
		\midrule
		Passaggio modulante & Quinte con alterazioni, sequenze reali & Transizione fluida & \ref{chap:progressioni_modulanti} \\
		\midrule
		Salti di quarta o quinta & Stati fondamentali (5) o catena di settime (7) & Chiarezza strutturale & \ref{chap:progressioni_armoniche}, \ref{chap:progressioni_settime} \\
		\midrule
		Modulo di due accordi ripetuto & Schemi I--IV e I--V, anche imitati & Sequenza regolare & \ref{chap:moduli}, \ref{chap:progressioni_settime} \\
		\midrule
		Salita cromatica & Progressione 5, 6/5 & Accumulo di tensione & \ref{chap:progressioni_cromatiche} \\
		\midrule
		Discesa cromatica & Progressione 2, 7, o basso cromatico discendente & Pathos, discesa controllata & \ref{chap:progressioni_cromatiche} \\
		\midrule
		Basso statico (nota pedale) & Variazione armonica sopra nota fissa & Stabilità con varietà & \ref{sec:pedale} \\
		\bottomrule
	\end{tabular}
\end{regola}
